\section*{Giratore Antoniou}

\noindent\colorbox{orange}{
  \begin{minipage}{1.0\linewidth}
    \subsection*{Richiesta (a)}
    Sulla base delle ipotesi di idealit\'a degli operazionali e di linearit\'a del circuito
    (quest'ultima verificata dalle misure precedenti), derivare un'espressione analitica
    dell'impedenza di ingresso del giratore (\textit{si suggerisce di procedere attraverso i seguenti
      passi, immaginando di scollegare la rete dal precedente circuito e di collegarla ad un
      generatore di tensione $V_{in}$ al suo ingresso A: in termini di $V_{in}$ determinare le correnti $I_4$
      ed $I_C$; quindi determinare $I_3$ ed $I_2$ considerando che i nodi B e C sono collegati agli
      ingressi dell'operazionale $U_2$; quindi determinare $I_1$ iterando lo stesso ragionamento per
      i nodi A e B collegati agli ingressi di $U_1$}).
  \end{minipage}
}
\subsection*{Risposta (a)}
Seguendo i suggerimenti, si calcola $I_4$, nell'ipotesi di idealit\'a dell'op-amp e operazione in regime
di linearit\'a, si applica il \textit{principio del corto-circuito virtuale}, per cui le tensioni ai nodi B e C
sono entrambi $V_{in}$ e da cui
\begin{equation*}
  I_4 = \frac{V_{in}}{R}
\end{equation*}
sempre per ipotesi di idealit\'a dell'op-amp gli ingressi invertenti non erogando corrente
\begin{equation*}
  I_3 = I_4
\end{equation*}
da cui si ottiene la tensione del nodo in comune con l'uscita dell'op-amp $U_1$ mediante la KCL
\begin{equation*}
  \begin{gathered}
    sC(V_{U_1} - V_{in}) = I_3 = I_4 = \frac{V_{in}}{R} \\ \implies
    V_{U_1} = \frac{1 + sRC}{sRC}V_{in}
  \end{gathered}
\end{equation*}
da cui si ottiene la corrente $I_3$
\begin{equation*}
  I_3 = \frac{V_{in} - V_{U_1}}{R}
\end{equation*}
sempre per ipotesi di idealit\'a dell'op-amp gli ingressi invertenti non erogando corrente
\begin{equation*}
  I_2 = I_3
\end{equation*}
da cui si ottiene la tensione del nodo in comune con l'uscita dell'op-amp $U_2$ mediante la KCL
\begin{equation*}
  \begin{gathered}
    \frac{V_{U_2} - V_{in}}{R} = I_2 = I_3 = \frac{V_{in} - V_{U_1}}{R} \\ \implies
    V_{U_2} = -\frac{1 - sRC}{sRC}V_{in}
  \end{gathered}
\end{equation*}
da cui si ottiene la corrente $I_1$, quindi finalmente anche l'impedenza d'ingresso $Z_{in}$
\begin{equation*}
  I_1 = \frac{V_{in} - V_{U_2}}{R} \implies \colorbox{blu}{\boxed{Z_{in} \equiv \frac{V_{in}}{I_1} = sR^2C}}
\end{equation*}

\noindent\colorbox{orange}{
  \begin{minipage}{1.0\linewidth}
    \subsection*{Richiesta (b)}
    Verificare che modulo ed argomento di detta impedenza siano in accordo con l'ipotesi di
    cui al punto 1d.
  \end{minipage}
}
\subsection*{Risposta (b)}
\begin{equation*}
  \begin{gathered}
    |Z_{in}| = R^2C \\
    \arg Z_{in} = \pi/2
  \end{gathered}
\end{equation*}

\noindent\colorbox{orange}{
  \begin{minipage}{1.0\linewidth}
    \subsection*{Richiesta (c)}
    Sostituendo il giratore con l'induttanza equivalente e con $C_1$ connessa, determinare la
    funzione di trasferimento del circuito $V_{OUT}/V_S$.
  \end{minipage}
}
\subsection*{Risposta (c)}
Sostituendo il contenuto del blocco con un induttore, in virt\'u della KCL
\begin{equation*}
  \frac{V_S - V_{OUT}}{R} = sCV_{OUT} + \frac{V_{OUT}}{sR^2C}
\end{equation*}
da cui si ottiene la funzione di trasferimento
\begin{equation*}
  \colorbox{blu}{\boxed{H(s) \equiv \frac{V_{OUT}}{V_S} = \frac{sRC}{1 + sRC + s^2R^2C^2}}}
\end{equation*}

\noindent\colorbox{orange}{
  \begin{minipage}{1.0\linewidth}
    \subsection*{Richiesta (d)}
    Dall'espressione ottenuta dedurre i valori attesi per i parametri del filtro (frequenza
    centrale, guadagno massimo e larghezza banda) e confrontarli con le misure di cui al
    punto 1c.
  \end{minipage}
}
\subsection*{Risposta (d)}
Restringendo la funzione di trasferimento all'asse immaginaria
\begin{equation*}
  A(j\omega) = \frac{jRC\omega}{(1 - R^2C^2\omega^2) + j(RC\omega)}
\end{equation*}
mediante la \ref{eq:da-tfunc-a-gain} e dividendo numeratore e denominatore per $RC\omega$ si ottiene
un'espressione per il guadagno che rende semplice l'identificazione del massimo
\begin{equation*}
  G(\omega) = \frac{1}{\sqrt{1 + \left(\frac{1 - R^2C^2\omega^2}{RC\omega}\right)^2}}
\end{equation*}
il massimo \'e un minimo del modulo del termine fra parentesi
\begin{equation*}
  \begin{gathered}
    \colorbox{blu}{\boxed{\Omega = \frac{1}{RC}}} \\
    \colorbox{blu}{\boxed{\max_{\omega}G(\omega) = G(\Omega) = 1}}
  \end{gathered}
\end{equation*}
si determina la frequenza di taglio
\begin{equation*}
  \begin{gathered}
    \frac{1}{1 + \left(\frac{1 - R^2C^2\omega^2}{RC\omega}\right)^2} = \frac12 \\ \implies
    R^2C^2\omega^2 \pm RC\omega - 1 = 0 \\ \implies
    \colorbox{blu}{\boxed{\omega = \frac{\sqrt5 \pm 1}{2RC}}}
  \end{gathered}
\end{equation*}
dunque la larghezza di banda
\begin{equation*}
  \colorbox{blu}{\boxed{\Delta\omega = \frac{1}{RC}}}
\end{equation*}

